\chapter*{}
%\thispagestyle{empty}
%\cleardoublepage

%\thispagestyle{empty}

\begin{titlepage}
 
 
\setlength{\centeroffset}{-0.5\oddsidemargin}
\addtolength{\centeroffset}{0.5\evensidemargin}
\thispagestyle{empty}

\noindent\hspace*{\centeroffset}\begin{minipage}{\textwidth}

\centering
%\includegraphics[width=0.9\textwidth]{imagenes/logo_ugr.jpg}\\[1.4cm]

%\textsc{ \Large PROYECTO FIN DE CARRERA\\[0.2cm]}
%\textsc{ INGENIERÍA EN INFORMÁTICA}\\[1cm]
% Upper part of the page
% 

 \vspace{3.3cm}

%si el proyecto tiene logo poner aquí
\includegraphics{imagenes/logo.png} 
 \vspace{0.5cm}

% Title

{\Huge\bfseries Título del proyecto\\
}
\noindent\rule[-1ex]{\textwidth}{3pt}\\[3.5ex]
{\large\bfseries Subtítulo del proyecto.\\[4cm]}
\end{minipage}

\vspace{2.5cm}
\noindent\hspace*{\centeroffset}\begin{minipage}{\textwidth}
\centering

\textbf{Autor}\\ {Nombre Apellido1 Apellido2 (alumno)}\\[2.5ex]
\textbf{Directores}\\
{Nombre Apellido1 Apellido2 (tutor1)\\
Nombre Apellido1 Apellido2 (tutor2)}\\[2cm]
%\includegraphics[width=0.15\textwidth]{imagenes/tstc.png}\\[0.1cm]
%\textsc{Departamento de Teoría de la Señal, Telemática y Comunicaciones}\\
%\textsc{---}\\
%Granada, mes de 201
\end{minipage}
%\addtolength{\textwidth}{\centeroffset}
\vspace{\stretch{2}}

 
\end{titlepage}






\cleardoublepage
\thispagestyle{empty}

\begin{center}
{\large\bfseries Estudio y Análisis de Metaheurísticas modernas para el problema de Selección de Características}\\
\end{center}
\begin{center}
MIGUEL GARCÍA LÓPEZ\\
\end{center}

%\vspace{0.7cm}
\noindent{\textbf{Palabras clave}: Metaheurística, selección, de características, binario, optimización, población, fitness, aprendizaje automático}\\

\vspace{0.7cm}
\noindent{\textbf{Resumen}}\\

En el ámbito del Machine Learning, algunos algoritmos, como KNN o SVM, muestran excelentes resultados, pero muchas veces estos requieren de un procesamiento previo para identificar las características más relevantes. Ya sea por escoger un problema amplio y complejo, por falta de conocimiento o por falta de contexto a la hora de recolectar los datos, es posible que se llegue a aglomerar información de más, dando lugar a conjuntos de datos inmensos. Incluso aquellos algoritmos que internamente realizan esta identificación de características pueden beneficiarse de un procesamiento adicional. Este preprocesamiento, conocido como selección de características (features selection), se considera un problema complejo de optimización combinatoria.\\[6pt]

Las metaheurísticas son algoritmos diseñados para resolver problemas de optimización complejos cuando los recursos son limitados. Aunque inicialmente se desarrollaron para abordar principalmente problemas combinatorios, en la actualidad se están proponiendo y aplicando cada vez más en problemas que implican variables continuas o reales. En respuesta a la creciente demanda en el campo de la selección de características, se han adaptado versiones especializadas de estas metaheurísticas para abordar este tipo de problemas combinatorios. Sin embargo, a pesar del número creciente de propuestas en este campo, las comparaciones objetivas entre ellas son limitadas. Aunque existen revisiones bibliográficas, muchas de ellas carecen de comparaciones adecuadas debido a la importancia y actualidad del problema de selección de características. Por lo tanto, hay una necesidad de estudios que proporcionen una evaluación comparativa más rigurosa y exhaustiva de las diferentes propuestas en este ámbito.\\[6pt]

En este trabajo de carácter científico, se llevará a cabo una revisión bibliográfica de diversas metaheurísticas recientes para abordar el problema de selección de características. Se estudiarán e implementarán aquellas consideradas más prometedoras, con el objetivo de construir un repertorio amplio y variado de propuestas. Posteriormente, se realizará un estudio comparativo exhaustivo utilizando diversos algoritmos de machine learning y conjuntos de datos representativos. Finalmente, se llevará a cabo un análisis crítico utilizando diversas métricas y valoraciones, como la tasa de acierto y el tiempo de ejecución, entre otras.
\cleardoublepage


\thispagestyle{empty}


\begin{center}
{\large\bfseries Study and Analysis of Modern Metaheuristics for the Feature Selection Problem}\\
\end{center}
\begin{center}
MIGUEL GARCÍA LÓPEZ\\
\end{center}

%\vspace{0.7cm}
\noindent{\textbf{Keywords}: Metaheuristic, feature selection, binary, optimization, population, fitness, machine learning}\\

\vspace{0.7cm}
\noindent{\textbf{Abstract}}\\

In the field of Machine Learning, some algorithms like KNN or SVM often yield excellent results, but they often require preprocessing to identify the most relevant features. Whether due to tackling a broad and complex problem, lack of domain knowledge, or collecting data without proper context, it's possible to gather excessive information, resulting in vast datasets. Even algorithms internally performing feature identification can benefit from additional processing. This preprocessing, known as feature selection, is considered a complex combinatorial optimization problem.\\[6pt]

Metaheuristics are algorithms designed to solve complex optimization problems when resources are limited. Initially developed mainly for combinatorial problems, they are increasingly proposed and applied to problems involving continuous or real variables. In response to the growing demand in feature selection, specialized versions of these metaheuristics have been adapted to tackle combinatorial problems. However, despite the increasing number of proposals in this field, objective comparisons between them are limited. Although there are literature reviews, many lack adequate comparisons due to the importance and timeliness of the feature selection problem. Therefore, there is a need for studies providing a more rigorous and exhaustive comparative evaluation of different proposals in this area.\\[6pt]

In this scientific work, a literature review of various recent metaheuristics for feature selection will be conducted. The most promising ones will be studied and implemented to build a broad and varied repertoire of proposals. Subsequently, a comprehensive comparative study will be conducted using various machine learning algorithms and representative datasets. Finally, a critical analysis will be carried out using various metrics and assessments, such as accuracy rate and execution time, among others.\\[6pt]

\chapter*{}
\thispagestyle{empty}

\noindent\rule[-1ex]{\textwidth}{2pt}\\[4.5ex]

Yo, \textbf{Miguel García López}, alumno de la titulación Ingeniería Informática de la \textbf{Escuela Técnica Superior
de Ingenierías Informática y de Telecomunicación de la Universidad de Granada}, con DNI 77159865E, autorizo la
ubicación de la siguiente copia de mi Trabajo Fin de Grado en la biblioteca del centro para que pueda ser
consultada por las personas que lo deseen.

\vspace{6cm}

\noindent Fdo: Miguel García López

\vspace{2cm}

\begin{flushright}
Granada a 29 de Enero de 2024.
\end{flushright}


\chapter*{}
\thispagestyle{empty}

\noindent\rule[-1ex]{\textwidth}{2pt}\\[4.5ex]

D. \textbf{Daniel Molina Cabrera}, Profesor del Departamento Ciencias de la Computación e Inteligencia Artificial de la Universidad de Granada.

\vspace{0.5cm}
\textbf{Informan:}

\vspace{0.5cm}

Que el presente trabajo, titulado \textit{\textbf{Estudio y Análisis de Metaheurísticas modernas para el problema de Selección de Características}},
ha sido realizado bajo su supervisión por \textbf{Miguel García López}, y autorizamos la defensa de dicho trabajo ante el tribunal
que corresponda.

\vspace{0.5cm}

Y para que conste, expiden y firman el presente informe en Granada a 29 de Enero de 2024.

\vspace{1cm}

\textbf{El director:}

\vspace{5cm}

\noindent \textbf{Daniel Molina Cabrera}

\chapter*{Agradecimientos}
\thispagestyle{empty}

       \vspace{1cm}


Poner aquí agradecimientos...

